% Parte «results» (clase section): la edita solo su dueño, declarado en roles.yaml.
% Escribe una oración por línea y separa los párrafos con una línea en blanco (02-convencion §4.3).
% El título de la sección y su etiqueta sec:results están en tex/outline/results.tex, del editor del artículo:
% aquí no van \section ni \chapter; usa \subsection y niveles menores (02-convencion REQ-CONV-94).

Esta parte analiza los resultados de los experimentos con sus figuras y tablas.
Las figuras van en \texttt{figures} y las tablas en \texttt{tables}, dentro de la carpeta de la parte.
